% !TEX root = ../../../main.tex
% Sources: LEC01, IMG_9111--IMG_9114.
\section{Rational gaps and irrational numbers}
\label{sec:analysis-i:rational-gaps}

\newthought{Counting is only one part of the story.} The rational numbers
are closed under the usual arithmetic operations, with division by zero
excluded. Yet even a simple geometric construction leads outside $\Q$.
The right triangle in \autoref{fig:unit-triangle} has legs of length one,
so its hypotenuse has square equal to two.
\begin{marginfigure}
  \centering
  \begin{tikzpicture}[analysis diagram,x=21mm,y=21mm]
  \draw (0,0) -- node[below] {$1$} (1,0)
    -- node[right] {$1$} (1,1)
    -- node[above left] {$\sqrt{2}$} cycle;
  \draw (0.87,0) -- (0.87,0.13) -- (1,0.13);
\end{tikzpicture}

  \caption[A unit right triangle.]{A unit right triangle produces the length $\sqrt{2}$.}
  \label{fig:unit-triangle}
\end{marginfigure}

\begin{proposition}
  The number $\sqrt{2}$ is irrational.
\end{proposition}
\begin{proof}
  Suppose $\sqrt{2}=p/q$, where $p,q\in\N$ have no common factor.
  Then $p^2=2q^2$, so $p^2$ is even. An odd integer has an odd square,
  hence $p=2r$ for some integer $r$. Substitution gives $q^2=2r^2$,
  so $q$ is even as well. This contradicts the choice of $p$ and $q$.
\end{proof}

The same divisibility argument extends to other primes and higher roots.
We take up that exercise in \autoref{sec:analysis-i:roots-and-powers},
after establishing existence of the roots themselves.

\subsection{A gap in the rational line}
We can describe the location of $\sqrt{2}$ without naming a rational number
at that location. Partition $\Q$ into
\begin{align*}
  L&=\{q\in\Q\mid q\leq0\text{ or }q^2<2\},\\
  U&=\{q\in\Q\mid q>0\text{ and }q^2>2\}.
\end{align*}
Both sets are nonempty, every rational belongs to one of them, and every
element of $L$ is less than every element of $U$. There is no rational
boundary between them. Indeed, if a rational boundary $c$ satisfied $c^2<2$,
a sufficiently small positive rational increment would leave its square
below two, contradicting its being a boundary. If $c^2>2$, a sufficiently
small rational decrement gives the opposite contradiction. The case $c^2=2$
has just been ruled out. Here a boundary must be positive, since $1\in L$.
The small increments can be checked directly from
$(c+h)^2=c^2+2ch+h^2$.

This is the idea behind a \emph{Dedekind cut}. Completing the rational line
means supplying its missing boundaries. An equivalent approach supplies
limits of rational Cauchy sequences. We will formulate both versions of
completeness in \autoref{sec:analysis-i:supremum}.
\marginnote{The condition $q\leq0$ in $L$ is necessary. The sets defined only
by $q^2<2$ and $q^2>2$ would put large negative rationals on the wrong side.}

\subsection{Limits can reveal further gaps}
The lecture also introduced $e$ as an irrational number obtained from
rational factorial partial sums. This is a different source of a gap from
the geometric example above. Its construction and irrationality proof
appear in \autoref{sec:analysis-i:denseness}, where monotone convergence
and geometric-tail estimates are available. For now it motivates the
question of which limits a number system contains.
